\subsection{Extended ablation study} 

This section contains the results of ablation experiments referenced from the main text (Section 4.2).

\subsubsection{Importance of positional encoding.}
There are two kinds of positional encoding in our model: spatial 
positional encoding and output positional encoding (object queries).
We experiment with various combinations of fixed and learned encodings for the spatial encodings. The positional encodings are shared between all layers,
and the output encoding (object queries) is always learned. The
results are given in Table~\ref{table:pos_enc}.
The first experiment removes both kinds of positional encoding altogether,
and, interestingly, the model still achieves 
more than 31 AP, losing 6 AP to the baseline.
Then, we add fixed sine spatial positional encoding and output encoding
to the input, as in the original transformer~\cite{Vaswani2017AttentionIA},
and find that this leads to 1.1 AP drop compared to passing positional encoding
at every attention layer. Learned spatial encoding at every attention layers gives
similar results to fixed sine encodings added to inputs.
Surprisingly, we find that removing the positional embeddings of the encoder
only leads to a minor AP drop.
In all cases, encodings are shared between all layers,
and output encoding (object queries) is always learned.
\setlength{\tabcolsep}{0.7em}
\begin{table}[h!]
    \centering\small
    \caption{Results for different positional encodings compared
    to baseline (last row), which has a fixed sine pos. encoding passed
    at every attention layer in both the encoder and decoder.
    Learned embedding is shared in all the layers.
    Not using positional encoding leads to a significant
    drop in AP. Interestingly, passing it in decoder only
    leads to a minor AP drop.\label{table:pos_enc}}
    \begin{tabular}{lcc|cc|cc}
        \toprule
        encoding & encoder & decoder & AP & $\Delta$ & \AP{50} & $\Delta$ \\
        \midrule
        none                & &     
        & 31.0 & -6.4 & 51.6 & -4.8 \\
        sine at input        & \chk & \chk
        & 36.2 & -1.1 & 55.1 & -1.3 \\
        learned at attention & \chk & \chk
        & 36.1 & -1.2 &	55.8 & -0.6 \\
        sine at attention & & \chk
        & 36.8 & -0.5 &	56.2 & -0.1 \\
        sine at attention & \chk & \chk
        & \textbf{37.3} & - & \textbf{56.4} & - \\
        \bottomrule
    \end{tabular}
\end{table}
Interestingly, the model without any positional encoding still achieves 
more than 31 AP, losing 6 AP to the baseline.
The model with absolute encoding added to the input only
loses 1.4 AP vs. baseline. 

\subsubsection{Cost ablations.} 
To evaluate the importance of different parts of the matching cost,
we train several models turning them on and off.
There are two main components to the cost: classification cost and box cost (\ref{eq:bloss}), consisting of GIoU and $\ell_1$ bounding box distance.
The classification cost is essential for training and cannot be turned off,
so we train a model without bounding box distance loss, and a model without the GIoU loss,
setting either $\lambda_{\rm iou}$ or $\lambda_{\rm L1}$ in the cost and in the loss to zero,
and compare with baseline.
Results are presented in Table~\ref{table:cost}.
\setlength{\tabcolsep}{0.7em}
\begin{table}[h!]
    \centering\small
    \caption{Effect of cost parts on AP.
    We train two models turning off box cost and GIoU cost,
    and observe that box cost gives better results on large objects,
    whereas GIoU cost is better on small.
    \detr (last row) combines both costs.
    }
    \begin{tabular}{ccc|cc|cc|ccc}
        \toprule
        class & $\ell_1$ & GIoU & AP & $\Delta$ & \AP{50} & $\Delta$ & \AP{S} & \AP{M} & \AP{L} \\
        \midrule
        \chk & \chk & &
        35.4 & -1.8 & 55.3 & -1.1 & 13.7 &	38.9 &	\textbf{56.4} \\
        \chk & & \chk &
        35.5 & -1.9 & 55.2 & -1.1 & 16.6 &	39.6 &	51.0 \\
        \chk & \chk & \chk &
        \textbf{37.3} & - & \textbf{56.4} & - & \textbf{17.8} &	\textbf{40.5} &	54.3 \\
        \bottomrule
    \end{tabular}
    \label{table:cost}
\end{table}
The GIoU cost shows better results on small objects (+2.9 AP difference),
whereas the bounding box cost improves results on large objects (+5.5 AP difference),
higher than the baseline. We only studied simple ablations of different costs (adding or removing one, and not retuning their weighting every time),
but perhaps other means of combining them could achieve different results.

\iffalse
\setlength{\tabcolsep}{0.7em}
\begin{table}[h!]
    \centering\small
    \begin{tabular}{ccc|cc|cc|ccc}
        \toprule
        class & bbox & GIoU & AP & $\Delta$ & \AP{50} & $\Delta$ & \AP{S} & \AP{M} & \AP{L} \\
        \midrule
        \chk & \chk & &
        35.4 & -1.8 & 55.3 & -1.1 & 13.7 &	38.9 &	\textbf{56.4} \\
        \chk & & \chk &
        35.5 & -1.9 & 55.2 & -1.1 & 16.6 &	39.6 &	51.0 \\
        \chk & \chk & \chk &
        \textbf{37.3} & - & \textbf{56.4} & - & \textbf{17.8} &	\textbf{40.5} &	54.3 \\
        \bottomrule
    \end{tabular}
    \caption{Effect of cost parts on AP.
    We train two models turning off box cost and GIoU cost,
    and observe that box cost gives better results on large objects,
    whereas GIoU cost is better on small.
    Our baseline (last row) combines both costs.
    }
    \label{table:cost}
\end{table}
\setlength{\tabcolsep}{0.5em}
\begin{figure}
\small
\begin{subfigure}[t!]{0.46\textwidth}
    \begin{tabular}{lccc}
        \toprule
        encoding & enc. & dec. & AP \\
        \midrule
        none                & &     
        & 31.0  \\
        sine @ input        & \chk & \chk
        & 36.2  \\
        learned @ attn & \chk & \chk
        & 36.1  \\
        sine @ attn & & \chk
        & 36.8  \\
        sine @ attn & \chk & \chk
        & \textbf{37.3} \\
        \bottomrule
    \end{tabular}
    \caption{Results for different positional encodings compared
    to baseline (last row), which has a fixed sine pos. encoding passed
    at every attention layer in both the encoder and decoder.
    \label{table:pos_enc}}
\end{subfigure}
\begin{subfigure}[t!]{0.52\textwidth}
    \begin{tabular}{cccccc}
        \toprule
        class & bbox & GIoU & AP %
        & \AP{S} %
        & \AP{L} \\
        \midrule
        \chk & \chk & &
        35.4 %
        & 13.7 %
        &	\textbf{56.4} \\
        \chk & & \chk &
        35.5 %
        & 16.6 %
        &	51.0 \\
        \chk & \chk & \chk &
        \textbf{37.3} %
        & \textbf{17.8} %
        &	54.3 \\
        \bottomrule
    \end{tabular}
    \caption{Effect of the parts of the loss on AP.
    We train two models turning off the box cost and the GIoU cost,
    and observe that the box cost gives better results on large objects,
    whereas GIoU cost is better on small.
    Our baseline (last row) combines both costs.
    \label{table:cost}}
\end{subfigure}
    \caption{TODO TODO TODO\label{fig:posencodingandcostablations}}
\end{figure}
\fi
